% TEMPLATE-MILC-v3.tex - plantilla maestra UNICA de las guias MILC v3.
%
% La usan las tres familias de guias, cada una con su driver:
%   · Grados 6-11  -> scripts/build-guias-g11.py
%   · Bebras       -> scripts/build-guias-bebras.py
%   · Semillero    -> scripts/build-guias-semillero.py
%
% Antes habia tres copias casi identicas (~400 lineas iguales, 6 distintas):
% cada mejora habia que aplicarla tres veces, y en la practica no se hacia
% (el motor de recursos vivia solo en dos de ellas). Lo que varia entre
% familias esta parametrizado en 6 placeholders que cada driver llena
% (se nombran aqui SIN su sintaxis real: el builder reemplaza por texto plano,
% tambien dentro de los comentarios, y un valor multilinea rompe el preambulo):
%
%   HEADER_IZQ        encabezado izquierdo de pagina
%   HEADER_DER        encabezado derecho de pagina
%   PORTADA_ETIQUETA  rotulo sobre el numero grande (GRADO/MOMENTO/MODULO)
%   PORTADA_SERIE     linea de serie de la portada
%   PORTADA_ID        identificador de la guia en la portada
%   PORTADA_CIRCULO   el circulito de grado (°), o vacio si no aplica
%
% El resto de claves las documenta content/guias/_SCHEMA.md. El script valida
% que no queden placeholders sin reemplazar antes de escribir el archivo final.

\documentclass[11pt,letterpaper]{article}
\usepackage[margin=1.75cm]{geometry}
\usepackage[table]{xcolor}
\usepackage{fontspec}
\usepackage[spanish]{babel}
\usepackage{tikz}
% Librerías TikZ para los diagramas de las guías (bloque `recursos:`):
%  · babel  → babel en español vuelve activos caracteres como «>», y TikZ los usa
%             en sus opciones (>=stealth, ->). Sin esta librería, cualquier
%             diagrama con flechas revienta con «\language@active@arg> has an
%             extra }». Debe ir ANTES de que se dibuje nada.
%  · positioning        → sintaxis «below left=2pt and 3pt».
%  · decorations.pathreplacing → llaves ({brace}) para señalar rangos.
\usetikzlibrary{babel,positioning,decorations.pathreplacing}
\usepackage{graphicx}
% Circuitos: el trabajo de electrónica se hace hoy con micro:bit y sus sensores
% integrados (MakeCode, sin cableado), así que no se necesitan esquemáticos.
% Si algún día entran montajes con protoboard, basta descomentar esta línea:
% los diagramas .tex/.tikz declarados en `recursos:` ya se incrustan solos.
% \usepackage{circuitikz}
\usepackage[most]{tcolorbox}
\usepackage{tabularx}
\usepackage{array}
\usepackage{fancyhdr}
\usepackage{titlesec}
\usepackage{ragged2e}
\usepackage{enumitem}
\usepackage{microtype}

% Helvetica Neue no trae la flecha U+2192, asi que cada «→» escrito literalmente
% en el YAML se compilaba a NADA, en silencio (462 flechas en 61 guias: rutas del
% tipo «paleta → tipografia → lockup» salian rotas). Mapeamos el caracter a su
% version matematica, que si tiene glifo. Asi el autor puede seguir escribiendo
% «→» comodamente en el YAML.
\usepackage{newunicodechar}
\newunicodechar{→}{\ensuremath{\rightarrow}}

\setmainfont{Helvetica Neue}
\setsansfont{Helvetica Neue}
\setlength{\parindent}{0pt}
\setlength{\parskip}{6pt}
\hyphenpenalty=200
\exhyphenpenalty=200
\pretolerance=400
\tolerance=2000
\setlength{\emergencystretch}{1em}
\renewcommand{\arraystretch}{1.25}
\setlist{leftmargin=5mm,itemsep=1mm,topsep=1mm}
\newcolumntype{Y}{>{\RaggedRight\arraybackslash}X}

\definecolor{milcMagenta}{HTML}{E6007E}
\definecolor{milcVino}{HTML}{5A0038}
\definecolor{milcTurquesa}{HTML}{0093A5}
\definecolor{milcVerde}{HTML}{58A923}
\definecolor{milcMostaza}{HTML}{E5B400}
\definecolor{milcOcre}{HTML}{B86600}
\definecolor{milcNegro}{HTML}{241816}
\definecolor{milcGris}{HTML}{F7F7F4}
\definecolor{milcLinea}{HTML}{E8E2DD}
\definecolor{milcGradePrimary}{HTML}{0066FF}
\definecolor{milcGradeDark}{HTML}{003D99}
\definecolor{milcGradeSoft}{HTML}{D6E8FF}
\definecolor{milcGradeAccent}{HTML}{CFE7FF}
\definecolor{milcGradeText}{HTML}{FFFFFF}
\color{milcNegro}

\pagestyle{fancy}
\fancyhf{}
\lhead{\small Tecnología e Informática · Grado 11}
\rhead{\small Guía 28 · MILC v3}
\cfoot{\small\thepage}
\renewcommand{\headrulewidth}{0pt}

\newtcolorbox{titlebox}[1]{
  enhanced,colback=#1,colframe=#1,arc=7mm,boxrule=0pt,
  left=7mm,right=7mm,top=3mm,bottom=3mm,
  before skip=8pt,after skip=8pt,
  fontupper=\sffamily\bfseries\Large\color{white}
}

\newtcolorbox{softbox}[3]{
  enhanced,colback=#2,colframe=#1,borderline west={4pt}{0pt}{#1},
  arc=2mm,boxrule=.4pt,left=4mm,right=4mm,top=3mm,bottom=3mm,
  before skip=6pt,after skip=8pt,title={#3},coltitle=white,
  colbacktitle=#1,fonttitle=\sffamily\bfseries,fontupper=\RaggedRight,
  attach boxed title to top left={xshift=3mm,yshift=-2mm},
  boxed title style={arc=2mm, boxrule=0pt}
}

\newtcolorbox{drawbox}[2]{
  enhanced,colback=white,colframe=#1,arc=4mm,boxrule=1pt,
  left=5mm,right=5mm,top=4mm,bottom=4mm,before skip=8pt,after skip=8pt,
  title={#2},coltitle=white,colbacktitle=#1,fonttitle=\sffamily\bfseries,
  fontupper=\RaggedRight,
  attach boxed title to top left={xshift=4mm,yshift=-2mm},
  boxed title style={arc=3mm, boxrule=0pt}
}

\newtcolorbox{infoband}[1]{
  enhanced,colback=#1!8,colframe=#1!50,arc=2mm,boxrule=.5pt,
  left=4mm,right=4mm,top=2mm,bottom=2mm,before skip=4pt,after skip=4pt,
  fontupper=\small\RaggedRight
}

% Puente narrativo entre secciones (contrato editorial Conéctate).
% Da continuidad: "Acabas de X. Ahora vas a Y porque Z."
% Visualmente: banda izquierda en color del grado, texto en itálica pequeña.
\newtcolorbox{puentebox}{
  enhanced,colback=milcGradeSoft,colframe=milcGradeDark,
  borderline west={3pt}{0pt}{milcGradeDark},
  arc=0mm,boxrule=0pt,left=5mm,right=4mm,top=2.5mm,bottom=2.5mm,
  before skip=4pt,after skip=8pt,
  fontupper=\itshape\small\color{milcNegro}\RaggedRight
}
% La flecha va en modo matematico a proposito: Helvetica Neue NO tiene el glifo
% U+2192, asi que \textrightarrow se compilaba a NADA («Missing character: There
% is no → in font Helvetica Neue») y el puente salia sin flecha. En modo
% matematico la flecha viene de la fuente matematica y si se dibuja.
\newcommand{\puente}[1]{%
  \begin{puentebox}%
    \textcolor{milcGradeDark}{$\rightarrow$}\hspace{2mm}#1%
  \end{puentebox}%
}

\newcommand{\linead}[1]{\noindent\color{milcLinea}\rule{#1}{0.5pt}\color{milcNegro}}
\newcommand{\respuesta}[1]{\par\vspace{2mm}\foreach \n in {1,...,#1}{\noindent\color{milcLinea}\rule{\linewidth}{0.5pt}\par\vspace{5mm}}\color{milcNegro}}
\newcommand{\checkbox}{\textcolor{milcGradeDark}{\fbox{\phantom{x}}}\hspace{5pt}}

\newcommand{\phasecard}[4]{
\begin{tcolorbox}[enhanced,colback=#3,colframe=#2,arc=3mm,boxrule=.5pt,height=3.05cm,valign=center,halign=center,left=1.5mm,right=1.5mm,top=2mm,bottom=2mm]
{\sffamily\bfseries\fontsize{22}{24}\selectfont\textcolor{#2}{#1}}\par
{\sffamily\bfseries\small #4}
\end{tcolorbox}
}

% ── Recursos visuales de la guía ──────────────────────────────────────
% Declarados en el bloque `recursos:` del YAML y emitidos por el builder.
% \guiaFigura   → imagen rasterizada o PDF (foto, carta celeste, montaje).
% \guiaDiagrama → fuente TikZ/circuitikz (.tex) incrustada como vector:
%                 es la vía para circuitos y esquemáticos editables.
% #1 = ruta relativa al .tex · #2 = pie (puede ir vacío)
\newcommand{\guiaFigura}[2]{%
\begin{center}
\includegraphics[width=0.88\linewidth,height=0.38\textheight,keepaspectratio]{#1}\\[1.5mm]
{\footnotesize\itshape\textcolor{milcNegro!75}{#2}}
\end{center}
}

\newcommand{\guiaDiagrama}[2]{%
\begin{center}
\input{#1}\\[1.5mm]
{\footnotesize\itshape\textcolor{milcNegro!75}{#2}}
\end{center}
}

\begin{document}
\thispagestyle{empty}
\begin{tikzpicture}[remember picture,overlay]
  \fill[white] (current page.south west) rectangle (current page.north east);
  \path[fill=milcGradePrimary,rounded corners=16mm]
    ([xshift=.55cm,yshift=-.55cm]current page.north west) rectangle
    ([xshift=-.55cm,yshift=3.65cm]current page.south east);

  \node[anchor=north west,text=milcGradeText] at ([xshift=2.0cm,yshift=-1.85cm]current page.north west)
    {\sffamily\bfseries\fontsize{14}{16}\selectfont GRADO};
  \node[anchor=north west,text=milcGradeText,opacity=.82] at ([xshift=2.0cm,yshift=-2.75cm]current page.north west)
    {\sffamily\bfseries\fontsize{9}{11}\selectfont SERIE GUÍAS · TECNOLOGÍA E INFORMÁTICA};

  \begin{scope}[shift={([xshift=-3.05cm,yshift=-2.55cm]current page.north east)}]
    \fill[white,opacity=.80] (-.62,-.52) rectangle (.62,.52);
    \draw[milcGradeText,line width=1pt] (-.62,-.52) rectangle (.62,.52);
    \fill[milcGradeDark] (-.42,-.38) rectangle (-.22,.06);
    \fill[milcGradeAccent] (-.10,-.38) rectangle (.10,.24);
    \fill[milcGradePrimary!60!black] (.22,-.38) rectangle (.42,.38);
  \end{scope}

  \node[anchor=west,text=milcGradeText] at ([xshift=1.9cm,yshift=-12.85cm]current page.north west)
    {\sffamily\bfseries\fontsize{124}{124}\selectfont 11};
  % El circulito de grado (°) solo aplica a las guias de grado («11°»); en
  % Bebras y Semillero el numero grande es un momento/modulo, no un grado.
  \draw[milcGradeText,line width=1.7pt]
    ([xshift=4.52cm,yshift=-11.68cm]current page.north west) circle (.13cm);

  \node[anchor=west,text=milcGradeText,text width=8.7cm,align=left] at ([xshift=10.35cm,yshift=-11.6cm]current page.north west)
    {
     {\sffamily\bfseries\fontsize{19}{22}\selectfont Ética y legal ---\\la palabra empeñada\\en lo digital}\\[4mm]
     {\sffamily\bfseries\fontsize{23}{26}\selectfont Guía 28-11-TIC}\\[2mm]
     {\sffamily\fontsize{13}{16}\selectfont Período 3 · Proyecto final emprendedor}\\[5mm]
     {\sffamily\bfseries\fontsize{11}{13}\selectfont Institución Educativa Sor María Juliana}\\[1mm]
     {\sffamily\fontsize{10}{12}\selectfont Autor: Dr. Álvaro Cárdenas Orozco}
    };

  \path[fill=milcGradeSoft,rounded corners=7mm]
    ([xshift=1.35cm,yshift=.85cm]current page.south west) rectangle
    ([xshift=-1.35cm,yshift=4.25cm]current page.south east);
  \draw[milcGradeDark,line width=.65pt,rounded corners=7mm]
    ([xshift=1.35cm,yshift=.85cm]current page.south west) rectangle
    ([xshift=-1.35cm,yshift=4.25cm]current page.south east);
  \node[anchor=center,text=milcNegro,text width=17.0cm,align=center] at ([yshift=2.55cm]current page.south)
    {\sffamily\fontsize{10}{12}\selectfont Metodología MILC · Apertura ancestral · Pensamiento computacional · Triángulo Dussel-Estoicismo-Floridi\\
    \textcolor{milcGradeDark}{\bfseries Producto:} 3 documentos legales y éticos del MVP en versión simple y honesta (a) Política de privacidad adaptada a la Ley 1581 de 2012 (habeas data colombiano), (b) Términos de uso del MVP con derechos y obligaciones del usuario, (c) Autoevaluación ética estructurada en las 4 preguntas de Floridi sobre el impacto del proyecto en la infosfera, con compromiso por escrito de qué se mantendrá, qué se ajustará y qué se descartará por razones éticas};
\end{tikzpicture}
\mbox{}
\newpage

\section*{Apertura: Ética y legal --- la palabra empeñada en lo digital}

\begin{softbox}{milcGradeDark}{milcGradeSoft}{Saber ancestral}
En los pueblos del Valle del Cauca y en las plazas de mercado de Cartago había una expresión que cerraba cualquier negocio sin necesidad de notario ni firma: \emph{``palabra empeñada''}. Cuando el tendero decía a la abuela \emph{``mañana le llega la panela''} y la abuela le daba el dinero por adelantado, no había contrato escrito: había \textbf{palabra empeñada}. Cuando el sastre prometía \emph{``el viernes le tengo el saco''} y aceptaba el adelanto, era palabra empeñada. Cuando el ebanista decía \emph{``este escaparate dura veinte años''}, era palabra empeñada. Romper la palabra no era falta legal: era \textbf{deshonor de oficio}, y producía consecuencias peores que cualquier multa, porque el ebanista que mentía perdía la clientela del barrio entero. La sabiduría ancestral tenía 3 reglas: (1) \textbf{nunca prometer lo que no se puede cumplir}; (2) \textbf{cumplir aunque cueste}; (3) \textbf{declarar a tiempo cuando algo va a fallar}. Esa triada se llamaba \emph{``código de honor del oficio''} y sostuvo economías enteras durante siglos antes de la era contractual.
\end{softbox}

\begin{softbox}{milcTurquesa}{milcGris}{Saber contemporáneo}
La \textbf{ética digital del emprendimiento} es exactamente la palabra empeñada del tendero traducida al ecosistema online. En Colombia, la \textbf{Ley 1581 de 2012} (Habeas Data) regula el tratamiento de datos personales: toda empresa o emprendimiento que recolecte datos de usuarios (nombre, correo, teléfono, ubicación, comportamiento) debe declarar explícitamente \textbf{qué datos recoge, para qué los usa, con quién los comparte, cuánto tiempo los guarda y cómo el usuario puede pedir su eliminación}. Esta declaración se llama \textbf{Política de Privacidad} y es obligatoria por ley para cualquier sitio que pida datos. Adicionalmente, los \textbf{Términos de Uso} establecen derechos y obligaciones del usuario y del prestador del servicio. Pero más allá de la ley está la \textbf{ética}: Luciano Floridi propone 4 preguntas que todo proyecto digital debería responder con honestidad sobre su impacto en la infosfera (el ecosistema de información compartido). La phronesis del emprendedor consiste en \textbf{tratar la ética como código de honor, no como trámite legal}.
\end{softbox}

\begin{softbox}{milcMostaza}{milcGris}{Pregunta-puente}
\textit{¿Qué sostenía la palabra empeñada del tendero antes de que existieran contratos digitales? ¿Y por qué la Política de Privacidad copiada-pegada de internet es la versión moderna del \emph{``firmar sin leer''} que los abuelos siempre advirtieron contra?}
\end{softbox}

\begin{softbox}{milcVerde}{milcGris}{Saber y saber hacer}
Al terminar podrás: (1) \textbf{identificar} qué datos recoge tu MVP, con qué propósito, con qué riesgos para el usuario, según la Ley 1581 de 2012; (2) \textbf{analizar} tu proyecto con las 4 preguntas éticas de Floridi sobre impacto en la infosfera (autenticidad, dignidad, privacidad, transparencia); (3) \textbf{evaluar} qué partes del MVP requieren ajuste ético antes de escalar; (4) \textbf{crear} 3 documentos en lenguaje claro: Política de Privacidad simple, Términos de Uso simples, Autoevaluación Ética con compromiso de ajustes.
\end{softbox}



\small
\begin{tabularx}{\textwidth}{>{\bfseries}p{3.0cm}Y}
\rowcolor{milcGradePrimary!12}Autor & Dr. Álvaro Cárdenas Orozco\\
DBA & Diseño, implementación y sustentación de un proyecto de impacto local (MEN, T\&I)\\
Referentes & Dussel (analéctica · sur global) · Lean Startup · Floridi · Estoicismo\\
Producto final & 3 documentos legales y éticos del MVP en versión simple y honesta (a) Política de privacidad adaptada a la Ley 1581 de 2012 (habeas data colombiano), (b) Términos de uso del MVP con derechos y obligaciones del usuario, (c) Autoevaluación ética estructurada en las 4 preguntas de Floridi sobre el impacto del proyecto en la infosfera, con compromiso por escrito de qué se mantendrá, qué se ajustará y qué se descartará por razones éticas\\
\end{tabularx}
\normalsize

\puente{Antes de empezar, mira el viaje completo. La sesión 7 cerró la dimensión financiera; hoy cierras la dimensión ético-legal. Las sesiones 9 (versión final del MVP) y 10 (sustentación pública) tomarán los compromisos éticos de hoy como parte del producto final. Sin política clara y autoevaluación honesta, el MVP escalado se convierte en pasivo legal y reputacional.}

\begin{titlebox}{milcGradeDark}
Ruta MILC: así vamos a aprender
\end{titlebox}
\begin{tabularx}{\textwidth}{>{\centering\arraybackslash}X>{\centering\arraybackslash}X>{\centering\arraybackslash}X>{\centering\arraybackslash}X}
\phasecard{1}{milcVerde}{milcGris}{Escucha\\[-1mm]\small Observo, escucho y pregunto.} &
\phasecard{2}{milcTurquesa}{milcGris}{Sistematización\\[-1mm]\small Pensamiento computacional.} &
\phasecard{3}{milcMagenta}{milcGris}{Praxis\\[-1mm]\small Construyo y aplico.} &
\phasecard{4}{milcVino}{milcGris}{Evaluación\\[-1mm]\small Reflexión liberadora.}
\end{tabularx}


\puente{Antes de teorizar, vas a hacer una \emph{auditoría de datos} de tu MVP: ¿qué información pides al usuario?, ¿dónde se almacena (Google Forms, Sheets, ManyChat, base propia)?, ¿quién tiene acceso?, ¿qué pasaría si esos datos se filtraran? Esa auditoría revela el alcance real de tu responsabilidad legal.}

\begin{titlebox}{milcVerde}
Fase 1 · Escucha
\end{titlebox}

\begin{softbox}{milcVerde}{milcGris}{Escena para escuchar}
\textbf{Actividad 1 · IDENTIFICA --- Auditoría de datos del MVP} (15 min · individual). Revisa tu MVP de la sesión 4 y haz una lista honesta de los datos que recoge del usuario. Reglas: (1) lista \textbf{cada dato} pedido (nombre, correo, teléfono, edad, ubicación, comportamiento de uso, dirección IP automática, foto, documento). (2) para cada dato, anota \textbf{dónde se almacena} (Google Sheets, ManyChat, Tally, base local, correo personal). (3) anota \textbf{quién tiene acceso} hoy (solo tú, herramienta sin tu control, terceros con permisos). (4) anota \textbf{cuánto tiempo} planeas guardarlo. (5) imagina el \textbf{escenario peor}: si esa base de datos se filtrara mañana en redes sociales, ¿qué daño concreto recibiría cada usuario?
\end{softbox}

\checkbox Listé cada dato que el MVP recoge.\\
\checkbox Para cada dato anoté dónde se almacena y quién tiene acceso.\\
\checkbox Definí tiempo de retención de cada dato.

\begin{infoband}{milcVerde}


\textbf{Lo que va al cuaderno (Actividad 1 · IDENTIFICA).} Título: «Actividad 1 --- Auditoría de datos». \textbf{Formato:} tabla 5+ filas × 5 columnas (Dato~|~Dónde se almacena~|~Quién accede~|~Cuánto tiempo se guarda~|~Daño en caso de filtración). \textbf{Sabes que terminaste cuando} podrías responder a un usuario preocupado con la verdad completa de qué hace tu MVP con sus datos.
\end{infoband}

\puente{Ya tienes el inventario de datos. Ahora vas a entender la \textbf{anatomía ético-legal} del MVP: los 5 elementos obligatorios de una Política de Privacidad bajo Ley 1581, los 6 elementos de unos Términos de Uso mínimos, las 4 preguntas de Floridi sobre la infosfera y los 5 errores típicos del emprendedor novato (copiar-pegar política ajena, omitir política, no leer lo que se firma, no permitir borrado de datos, no avisar al usuario de cambios).}

\begin{titlebox}{milcTurquesa}
Fase 2 · Sistematización con pensamiento computacional
\end{titlebox}

\begin{softbox}{milcTurquesa}{milcGris}{Cuatro pilares aplicados al tema}
Una política de privacidad copiada de internet sin entenderla es la versión digital del contrato firmado sin leer: protege legalmente al que copió, no al usuario. Para que tu proyecto cierre con dignidad ético-legal necesitas la \textbf{anatomía de los 3 documentos} (privacidad, términos, autoevaluación) y las \textbf{4 preguntas éticas} de Floridi aplicadas con honestidad.
\end{softbox}

\small
\begin{tabularx}{\textwidth}{>{\bfseries\RaggedRight\arraybackslash}p{3.6cm}Y}
\rowcolor{milcTurquesa!90}\color{white}Pilar & \color{white}\bfseries Aplicación al tema\\
1. Descomposición & La \textbf{Política de Privacidad} bajo Ley 1581 de 2012 tiene 5 elementos obligatorios: (1) \textbf{Qué datos se recogen}: lista explícita (no \emph{``datos personales''}, sino \emph{``nombre, correo electrónico, número de WhatsApp''}). (2) \textbf{Para qué se usan}: finalidad concreta (\emph{``contactar al usuario para coordinar la demo''}). (3) \textbf{Con quién se comparten}: terceros nombrados (\emph{``Google Workspace como proveedor de Sheets, ManyChat como proveedor de chatbot''}). (4) \textbf{Cuánto tiempo se guardan}: plazo concreto en meses o años. (5) \textbf{Derechos del usuario}: cómo pedir acceso, corrección, eliminación de sus datos. Si tu política no tiene estos 5, no es política: es texto decorativo.\\
2. Reconocimiento de patrones & Los \textbf{Términos de Uso} mínimos tienen 6 elementos: (1) Quién opera el servicio (nombre del proyecto + responsable). (2) Qué ofrece el servicio (descripción honesta del MVP). (3) Qué se compromete a hacer (cumplir lo prometido, responder en X tiempo). (4) Qué espera del usuario (uso adecuado, no compartir cuenta, no usar para X). (5) Qué pasa si algo falla (responsabilidades limitadas, garantías). (6) Cómo se contacta al responsable (correo o teléfono real). En lenguaje claro, no jurídico copiado.\\
3. Abstracción & Las \textbf{4 preguntas éticas de Floridi} sobre impacto en la infosfera: (1) \textbf{Autenticidad}: ¿mi MVP fomenta información verídica o ruido informacional? (2) \textbf{Dignidad}: ¿trata al usuario como sujeto con autonomía o como objeto a manipular? (3) \textbf{Privacidad}: ¿recoge solo lo necesario o acumula datos por si acaso? (4) \textbf{Transparencia}: ¿el usuario entiende lo que ocurre o el sistema es opaco a propósito? La phronesis ética consiste en \textbf{responder con honestidad aunque la respuesta exija ajustar el proyecto}.\\
4. Algoritmo conceptual & Algoritmo: (1) escribe la Política de Privacidad con los 5 elementos en lenguaje claro; (2) escribe Términos de Uso con los 6 elementos; (3) aplica las 4 preguntas de Floridi a tu MVP con respuesta honesta a cada una; (4) identifica al menos 1 ajuste concreto que harás por razones éticas (\emph{``no pediré la edad porque no la necesito''}, \emph{``borraré los datos de inscripciones no convertidas a los 30 días''}); (5) integra los 3 documentos como link visible desde la captura del MVP.\\
\end{tabularx}
\normalsize

\begin{drawbox}{milcTurquesa}{Anatomía ético-legal del MVP}
\textbf{Política de Privacidad:} 5 elementos (qué datos, para qué, con quién, cuánto tiempo, derechos). \\[1mm] \textbf{Términos de Uso:} 6 elementos (quién, qué, compromiso, expectativa, falla, contacto). \\[1mm] \textbf{4 preguntas Floridi:} autenticidad, dignidad, privacidad, transparencia. \\[1mm] \textbf{Compromiso de ajuste:} al menos 1 cambio por razones éticas.
\end{drawbox}

\begin{softbox}{milcMostaza}{milcGris}{Errores comunes y su corrección}
(1) \textbf{Copiar-pegar política ajena} sin adaptar a los datos reales del MVP. (2) \textbf{Omitir política}: arrancar sin política equivale a operar sin licencia de oficio. (3) \textbf{Política en lenguaje jurídico opaco}: el usuario no la entiende y por tanto no la acepta realmente. (4) \textbf{No permitir borrado de datos}: viola el derecho de habeas data. (5) \textbf{No avisar al usuario de cambios}: cualquier cambio en la política exige aviso. (6) \textbf{Confundir ética con legal}: cumplir la ley es mínimo, no máximo. (7) \textbf{Autoevaluación ética con respuestas decorativas}: \emph{``mi proyecto respeta al usuario''} sin nada concreto que lo demuestre.
\end{softbox}

\begin{infoband}{milcTurquesa}


\textbf{Lo que va al cuaderno (Actividad 2 · ANALIZA + EVALÚA).} Título: «Actividad 2 --- Anatomía ético-legal». \textbf{Formato:} (a) lista de los 5 elementos de la Política con su contenido para tu MVP; (b) lista de los 6 elementos de los Términos; (c) las 4 preguntas de Floridi respondidas con honestidad. \textbf{Sabes que terminaste cuando} podrías responder a una pregunta dura de un usuario sobre sus datos sin dudar.
\end{infoband}

\puente{Tienes el método. Toca redactar. Vas a escribir los 3 documentos en lenguaje claro (no jurídico copiado de internet), aplicar las 4 preguntas de Floridi a tu proyecto con honestidad y declarar qué vas a mantener, ajustar o descartar por razones éticas.}

\begin{titlebox}{milcMagenta}
Fase 3 · Praxis con pensamiento computacional
\end{titlebox}

\begin{softbox}{milcMagenta}{milcGris}{Construyendo el producto con los cuatro pilares}
\textbf{Actividad 3 · CREA --- 3 documentos ético-legales del MVP} (40 min · individual). Hora de redactar. Escribes los 3 documentos en lenguaje claro y honesto, identificas los ajustes éticos que harás antes de la sustentación, e integras los documentos al MVP como links visibles.
\end{softbox}

\small
\begin{tabularx}{\textwidth}{>{\bfseries\RaggedRight\arraybackslash}p{3.6cm}Y}
\rowcolor{milcMagenta!90}\color{white}Pilar & \color{white}\bfseries Aplicación al producto\\
Descomposición & Tus 3 documentos se construyen siguiendo plantillas mínimas: (a) \textbf{Política de Privacidad}: 1 página máximo, lenguaje claro, los 5 elementos obligatorios estructurados como respuestas a 5 preguntas (\emph{``¿qué datos recogemos?''}, \emph{``¿para qué los usamos?''}, etc.). (b) \textbf{Términos de Uso}: 1 página máximo, los 6 elementos en lenguaje cotidiano. (c) \textbf{Autoevaluación Ética}: 4 preguntas de Floridi con respuesta honesta de 3-4 líneas cada una + compromisos de ajuste concretos.\\
Patrones & Sigue el patrón \textbf{``lenguaje cotidiano sobre jerga legal''}: tus documentos no son contratos formales, son acuerdos honestos entre tú y tu usuario. Si el lenguaje suena a abogado de televisión, está mal escrito. Mejor frases cortas, segunda persona (\emph{``recogemos tu correo''}, \emph{``puedes pedir borrar tus datos en cualquier momento''}). La phronesis legal consiste en \textbf{que el usuario común entienda lo que firma}.\\
Abstracción & El compromiso de ajuste es la pieza más importante de la autoevaluación: no basta con responder las preguntas; debes declarar al menos 1 acción concreta. Ejemplos: (a) \emph{``Reduciré los campos del formulario de 8 a 4 porque no necesito edad, ciudad, ocupación, ni teléfono''}. (b) \emph{``Cambiaré el botón de \emph{Aceptar} por uno que diga \emph{Entiendo y acepto cómo se usan mis datos}}. (c) \emph{``Eliminaré la base de inscripciones no convertidas cada 30 días en lugar de retenerla indefinidamente''}. Sin compromiso de ajuste concreto, la autoevaluación es decorativa.\\
Algoritmo de armado & Algoritmo: (1) redacta Política en lenguaje claro; (2) redacta Términos en lenguaje claro; (3) responde las 4 preguntas de Floridi con honestidad; (4) declara compromiso de al menos 1 ajuste; (5) integra los 3 documentos al MVP como link visible desde la captura; (6) pídele a un familiar que no conoce tu proyecto que los lea y te diga si entendió.\\
\end{tabularx}
\normalsize

\begin{drawbox}{milcMagenta}{Checklist antes de cerrar la sesión}
\checkbox Política de Privacidad con los 5 elementos en lenguaje claro. \\[1mm] \checkbox Términos de Uso con los 6 elementos en lenguaje claro. \\[1mm] \checkbox Las 4 preguntas de Floridi respondidas con honestidad. \\[1mm] \checkbox Al menos 1 compromiso de ajuste ético escrito. \\[1mm] \checkbox 3 documentos integrados al MVP como link visible. \\[1mm] \checkbox Un familiar leyó los documentos y los entendió. \\[1mm] \checkbox Los ajustes se aplicarán antes de la sesión 9.
\end{drawbox}

\begin{softbox}{milcMagenta}{milcGris}{Plantilla de trabajo}
\textbf{Política de Privacidad.} \textit{Qué datos + para qué + con quién + cuánto tiempo + derechos del usuario.} \\[1mm] \textbf{Términos de Uso.} \textit{Quién opera + qué ofrece + qué se compromete + qué espera + responsabilidad + contacto.} \\[1mm] \textbf{Autoevaluación.} \textit{Autenticidad / Dignidad / Privacidad / Transparencia: respuesta honesta + compromiso de ajuste.}
\end{softbox}

\begin{infoband}{milcMagenta}


\textbf{Lo que va al cuaderno (Actividad 3 · CREA).} Título: «Actividad 3 --- Mis 3 documentos ético-legales». \textbf{Formato:} (a) Política de Privacidad completa (1 página); (b) Términos de Uso completos (1 página); (c) Autoevaluación con las 4 preguntas respondidas + compromiso de ajuste. \textbf{Extensión:} 3 páginas total. \textbf{Sabes que terminaste cuando} se lo podrías mostrar a tu profesor de filosofía o a un asesor legal y defenderlos con coherencia.
\end{infoband}

\puente{Lo que sigue resume \emph{qué} entregas hoy: Política de Privacidad simple + Términos de Uso simples + Autoevaluación Ética con compromiso. El criterio principal: que un usuario común leyendo tus documentos pueda entender en 5 minutos qué pasa con sus datos y qué se compromete contigo, sin necesidad de abogado.}

\begin{titlebox}{milcMostaza}
Producto de la sesión
\end{titlebox}
\textbf{Política + Términos + Autoevaluación ética con compromiso de ajuste}.
\begin{softbox}{milcMostaza}{milcGris}{Criterios de calidad}
(1) Política de Privacidad con los 5 elementos de Ley 1581 en lenguaje claro. (2) Términos de Uso con los 6 elementos mínimos. (3) Las 4 preguntas de Floridi respondidas con honestidad. (4) Al menos 1 compromiso de ajuste concreto. (5) Documentos integrados al MVP como link visible. (6) Validación con lectura por persona común no familiarizada.
\end{softbox}

\puente{Antes de cerrar, mira la dimensión ética desde las cinco dimensiones humanas. Empeñar la palabra digital es honor del oficio; copiar políticas sin leer es deshonor escondido tras lenguaje legal.}

\begin{titlebox}{milcVino}
Fase 4 · Evaluación liberadora — cinco dimensiones
\end{titlebox}

\begin{softbox}{milcVino}{milcGris}{Sentido de la evaluación}
Evaluar no es solo revisar una nota. Es reconocer qué aprendí, qué cambió en mí, qué emociones manejé, cómo me relaciono con la ciudad y la comunidad, y qué le dejaré a quien viene detrás.
\end{softbox}

\textbf{Mi reflexión liberadora en cinco dimensiones.} Completa cada frase iniciada:

\small
\begin{tabularx}{\textwidth}{>{\bfseries\RaggedRight\arraybackslash}p{3.4cm}Y>{\RaggedRight\arraybackslash}p{4.6cm}}
\rowcolor{milcVino}\color{white}Dimensión & \color{white}\bfseries Mi respuesta (frame starter) & \color{white}\bfseries Algo que descubrí\\
1. Personal & Lo que más cambió en mí fue \linead{6.5cm} & \linead{4.2cm}\\
2. Emocional & La emoción más fuerte fue \linead{2.5cm} y la regulé \linead{3cm} & \linead{4.2cm}\\
3. Ciudadana & En democracia esto importa porque \linead{6cm} & \linead{4.2cm}\\
4. Local (Cartago) & En mi barrio sería útil para \linead{6.5cm} & \linead{4.2cm}\\
5. Intergeneracional & A mi abuela/abuelo le diría \linead{6.5cm} & \linead{4.2cm}\\
\end{tabularx}
\normalsize

\begin{infoband}{milcVino}
\textbf{Para la próxima sustentación.} Vuelve a esta tabla en seis meses. Verás que las respuestas cambian — no porque lo aprendido cambie, sino porque tú cambias. La evaluación liberadora es una conversación contigo a través del tiempo.
\end{infoband}


\puente{Y para terminar, tres voces sobre la palabra y el honor. Dussel sobre el derecho del afectado a ser informado y consultado, Epicteto sobre la coherencia entre palabra y acción, Floridi sobre la infosfera como entorno ético compartido.}

\begin{titlebox}{milcGradeDark}
Triángulo de pensamiento · cierre filosófico
\end{titlebox}

\begin{softbox}{milcVino}{milcGris}{Dussel · lente del nosotros}
\textit{``El derecho del afectado a ser informado y consultado es la primera condición de la ética profesional.''}

Tu usuario es el afectado de tu proyecto: sus datos quedan en tu sistema, su tiempo se invierte en tu MVP, su confianza sostiene tu modelo. La filosofía de la liberación pide reconocerlo como sujeto con derechos, no como objeto del que se extrae información. La Política de Privacidad y los Términos no son trámite: son el acto formal de tratarlo como sujeto. Quien escribe política copiada sin leerla repite la lógica colonial del contrato impuesto al pueblo.

\textbf{¿Mis documentos legales reconocen al usuario como sujeto informado, o lo tratan como objeto que firma sin entender?}
\end{softbox}
\linead{\linewidth}\\[1mm]
\linead{\linewidth}

\begin{softbox}{milcMostaza}{milcGris}{Estoicismo · Epicteto · cuidado interior}
\textit{``Lo prometido y lo hecho deben coincidir; la grieta entre ambos es la medida del defecto del carácter.''}

Es tentador escribir política bonita y operar distinto. La disciplina estoica recomienda lo opuesto: \emph{lo prometido y lo hecho deben coincidir}. Si tu política dice que borras datos al mes y los guardas un año, no tienes problema legal solo: tienes problema de carácter profesional. La phronesis del oficio digital consiste en \textbf{escribir solo lo que vas a cumplir y cumplir todo lo que escribes}.

\textbf{¿Lo que prometo en mis documentos coincide con lo que mi MVP realmente hace, o hay grieta entre palabra y acción?}
\end{softbox}
\linead{\linewidth}\\[1mm]
\linead{\linewidth}

\begin{softbox}{milcTurquesa}{milcGris}{Floridi · lente de la infoesfera}
\textit{``La infosfera es el entorno ético compartido donde toda acción digital tiene consecuencias para los demás.''}

Tu MVP no opera en vacío: vive en la \emph{infosfera}, el ecosistema de información que todos compartimos. Cada decisión sobre datos, transparencia, autenticidad afecta a otros usuarios e indirectamente a la confianza colectiva en el oficio digital. Las 4 preguntas de Floridi te entrenan en pensar tu proyecto como \emph{ciudadano de la infosfera}, no como operador aislado. Quien se entrena hoy en esa ética escala con dignidad; quien la posterga acumula deuda ética con la sociedad.

\textbf{¿Mi MVP aporta autenticidad, dignidad, privacidad y transparencia a la infosfera, o las erosiona aunque sea un poco?}
\end{softbox}
\linead{\linewidth}\\[1mm]
\linead{\linewidth}

\begin{titlebox}{milcVino}
Compromiso final
\end{titlebox}

\small
\begin{tabularx}{\textwidth}{>{\RaggedRight\arraybackslash}p{3.0cm}YYY}
\rowcolor{milcVino}\color{white}\bfseries Criterio & \color{white}\bfseries Lo logré & \color{white}\bfseries En proceso & \color{white}\bfseries Necesito apoyo\\
Comprensión del tema & Explico ideas centrales con mis palabras. & Reconozco ideas pero falta precisión. & Necesito repasar conceptos base.\\
Pensamiento computacional & Apliqué los 4 pilares al diseñar mi producto. & Apliqué algunos pilares. & Necesito releer la fase 2.\\
Aplicación y producto & Mi producto cumple los criterios y se puede revisar. & Producto funcional con ajustes pendientes. & Aún en construcción.\\
Reflexión 5 dimensiones & Respondí cada dimensión con honestidad. & Respondí algunas con detalle. & Mis respuestas son muy generales.\\
Triángulo de pensamiento & Conecté las 3 voces con mi trabajo real. & Conecté al menos una. & No relaciono las voces con mi proceso.\\
\end{tabularx}
\normalsize

\textbf{Hoy entendí que:}\\[1mm]
\linead{\linewidth}\\[1mm]
\linead{\linewidth}

\textbf{Mi compromiso para la próxima sesión es:}\\[1mm]
\linead{\linewidth}\\[1mm]
\linead{\linewidth}

\begin{softbox}{milcMostaza}{milcGris}{Cita de despedida · Marco Aurelio}
\textit{``El obstáculo en el camino se vuelve el camino. Lo que estorba al camino se vuelve camino.''}\\[1mm]
Cada error de hoy, bien observado, se convierte en pieza del camino que pisarás mañana.
\end{softbox}

\small
\begin{tabularx}{\textwidth}{>{\bfseries}p{3.6cm}Y}
\rowcolor{milcGradeDark!12}Nombre completo: & \linead{10cm}\\
Fecha: & \linead{4cm} \quad Curso/grupo: \linead{4cm}\\
Mi firma: & \linead{9cm}\\
\end{tabularx}
\normalsize

\begin{infoband}{milcGradeDark}
\textbf{Para la próxima sesión.} Llega con los 3 documentos integrados al MVP y los ajustes éticos aplicados. En la sesión 9 vas a producir la \textbf{versión final del MVP}: el MVP pulido con todos los hallazgos del periodo (datos S5, presupuesto S7, ética S8). Sin los ajustes hechos esta semana, la versión final llegará con deuda ética visible.
\end{infoband}

\end{document}
